\par\medskip
\noindent\begin{minipage}{\textwidth}
\makeatletter\def\@captype{figure}\makeatother
\centering
\begin{minipage}[t]{0.53\textwidth}
\centering
\begin{tikzpicture}[scale=1.15, line cap=round, line join=round]
  \coordinate (v0) at (-2,1.05);
  \coordinate (v1) at (-2.95,-0.65);
  \coordinate (v2) at (-1.05,-0.65);
  \draw[thick,fill=blue!3] (v0)--(v1)--(v2)--cycle;
  \node[above] at (v0) {\(0\)};
  \node[below left] at (v1) {\(1\)};
  \node[below right] at (v2) {\(2\)};
  \node at (-2,-0.12) {\(\sigma\)};
  \node[font=\footnotesize,align=center] at (-2,-1.42)
    {one barycentric\\chamber \(\sigma\simeq\Delta^2\)};
  \draw[-{Stealth[length=2mm]},thick] (-0.25,0.12)--(-0.90,0.12)
    node[midway,above] {\(X\)};
  \begin{scope}[shift={(1,0)}]
    \path[fill=gray!22,draw=black,thick]
      (30:1.40)--(90:1.40)--(150:1.40)--(210:1.40)--
      (270:1.40)--(330:1.40)--cycle;
    \path[fill=white,draw=none]
      (30:0.76)--(90:0.76)--(150:0.76)--(210:0.76)--
      (270:0.76)--(330:0.76)--cycle;
    \node[font=\small] at (60:0.97) {\(\{0\}\)};
    \node[font=\small] at (120:0.97) {\(\{0,1\}\)};
    \node[font=\small] at (180:0.97) {\(\{1\}\)};
    \node[font=\small] at (240:0.97) {\(\{1,2\}\)};
    \node[font=\small] at (300:0.97) {\(\{2\}\)};
    \node[font=\small] at (0:0.87) {\(\{0,2\}\)};
    \node[font=\small] at (0,0) {tile};
  \end{scope}
  \node[font=\small,align=center] at (-0.55,-2.15)
    {(a) A chamber and its resolved tile};
\end{tikzpicture}
\end{minipage}
\hfill
\begin{minipage}[t]{0.43\textwidth}
\centering
\begin{tikzpicture}[scale=1.15, line cap=round, line join=round]
  \fill[blue!12] (0,0) rectangle (1.25,1.25);
  \fill[orange!15] (-1.25,0) rectangle (0,1.25);
  \fill[orange!15] (0,-1.25) rectangle (1.25,0);
  \fill[blue!12] (-1.25,-1.25) rectangle (0,0);
  \fill[gray!25] (-0.10,-1.25) rectangle (0.10,1.25);
  \fill[gray!25] (-1.25,-0.10) rectangle (1.25,0.10);
  \draw[thick] (-1.25,-1.25) rectangle (1.25,1.25);
  \draw[-{Stealth[length=1.6mm]}] (-1.4,0)--(1.5,0)
    node[right,font=\small] {\(Z_S\)};
  \draw[-{Stealth[length=1.6mm]}] (0,-1.4)--(0,1.5)
    node[above,font=\small] {\(Z_T\)};
  \node at (0.68,0.67) {\(\zeta\)};
  \node at (-0.68,0.67) {\(s_S\zeta\)};
  \node at (0.68,-0.67) {\(s_T\zeta\)};
  \node at (-0.68,-0.67) {\(s_Ss_T\zeta\)};
  \node[font=\small] at (1.03,1.04) {\(+\)};
  \node[font=\small] at (-1.03,1.04) {\(-\)};
  \node[font=\small] at (1.03,-1.04) {\(-\)};
  \node[font=\small] at (-1.03,-1.04) {\(+\)};
  \node[font=\small,align=center] at (0,-2.15)
    {(b) \(u_S=Z_S^2,\quad u_T=Z_T^2\)};
\end{tikzpicture}
\end{minipage}
\caption{The reflected resolution in dimension two, following
Gaifullin's permutohedral construction \cite{Gaifullin2008}. In (a), one
barycentric chamber \(\sigma\simeq\Delta^2\) is paired with its resolved
tile. The six facets are indexed by the nonempty proper subsets of
\(\{0,1,2\}\); adjacent labels are nested. The map \(X\) sends the tile
interior diffeomorphically to the chamber interior. In (b), at a corner
with labels \(S\subset T\), a full state \(\zeta=(\sigma,p,t)\) and its three reflected states occupy the four
quadrants of signed coordinates. The commutation
\(s_Ss_T=s_Ts_S\) and the independent bit toggles give four distinct
quadrants of one smooth chart. The displayed signs compare each chamber
sign with that of \(\zeta\): blue agrees and orange is opposite. The gray collar
and strips represent the same seam schematically; their widths are
illustrative, and \(X\) is not affine.}
\label{fig:resolution}
\end{minipage}
\par\medskip
